\chapter{Borderline Personality}
\index{Borderline Personality}

\section*{Definition and Overview}
Borderline personality refers to a pattern of unstable emotions, relationships, self-image, and behaviour that is captured more fully in the diagnosis of borderline personality disorder (BPD). It is sometimes called emotionally unstable personality disorder (EUPD). The core feature is marked emotional dysregulation: intense feelings that shift rapidly, a chronic sense of emptiness, fear of abandonment, impulsive behaviour, and unstable but intense relationships with others. People with BPD may also experience paranoia or dissociation under stress and, in a substantial proportion, self-harm and suicidal thoughts. The term "personality disorder" does not mean the person is difficult or untreatable; it describes a long-standing pattern of inner experience that deviates markedly from cultural expectations and causes distress and impairment. With appropriate treatment, most people with borderline personality improve substantially.

\section*{Epidemiology}
Borderline personality disorder is estimated to affect around 1--2\% of the general population. It is more commonly diagnosed in women, although community studies suggest the sex difference may be smaller than clinic samples imply. It is highly prevalent in clinical settings, being diagnosed in a significant proportion of psychiatric inpatients and outpatients, and it is the personality disorder most strongly associated with self-harm and suicide attempts. Onset is typically in adolescence or early adulthood, and the pattern usually persists across the adult years, although the intensity of symptoms tends to diminish with age. It is frequently associated with other mental health conditions, including depression, anxiety, substance use, and eating disorders.

\section*{Aetiology and Pathophysiology}
\begin{itemize}
    \item \textbf{Genetics:} twin and family studies show a heritable component, with a family history of BPD or related traits increasing risk.
    \item \textbf{Childhood adversity:} emotional, physical, or sexual abuse, neglect, and early separation from caregivers are strongly associated with BPD and are among the most robust risk factors.
    \item \textbf{Attachment:} disrupted early attachment can impair the development of the skills needed to soothe strong emotions and to trust others.
    \item \textbf{Neurobiology:} imaging and physiological studies show heightened reactivity of the amygdala to emotional stimuli and impaired functioning of the prefrontal cortex, which normally down-regulates emotional responses.
    \item \textbf{Diminished capacity model:} early invalidation teaches a person that their feelings are wrong or unimportant, so emotions are experienced as overwhelming and difficult to label or manage.
    \item \textbf{Personality traits:} high emotional sensitivity, high reactivity, and a slow return to baseline after an emotional trigger combine to produce the characteristic instability.
    \item \textbf{Stress sensitivity:} people with BPD are more sensitive to perceived rejection, and even small interpersonal events can trigger intense distress and impulsive responding.
\end{itemize}

\section*{Clinical Features}
\begin{itemize}
    \item Intense, rapidly shifting emotions, with anger, sadness, and anxiety often triggered by small events
    \item A chronic feeling of emptiness and an unstable sense of self and identity
    \item Intense but unstable relationships, oscillating between idealisation and devaluation of others
    \item A frantic fear of abandonment and desperate efforts to avoid being left alone
    \item Impulsive behaviour in spending, substance use, sex, or driving
    \item Recurrent self-harm, such as cutting or burning, and suicidal threats or attempts
    \item Episodes of paranoia or dissociation (feeling unreal or detached) under stress
    \item Mood swings that resolve within hours or days, rather than the sustained episodes seen in mood disorders
\end{itemize}

\section*{Differential Diagnosis}
\begin{itemize}
    \item \textbf{Bipolar disorder:} distinguished by sustained episodes of elevated or depressed mood lasting days to weeks, rather than brief emotional swings.
    \item \textbf{Depression and anxiety disorders:} chronic low mood or anxiety can coexist with, or mimic, the emotional instability of BPD.
    \item \textbf{Post-traumatic stress disorder:} can coexist with BPD; both follow trauma, but PTSD is defined by re-experiencing and avoidance of the trauma.
    \item \textbf{Other personality disorders:} histrionic, narcissistic, and dependent traits overlap with some features of BPD.
    \item \textbf{Substance use disorders:} intoxication and withdrawal can produce impulsivity and mood instability that resolve with abstinence.
    \item \textbf{Attention deficit hyperactivity disorder:} impulsivity and emotional reactivity in adulthood can be confused with BPD.
    \item \textbf{Dissociative disorders:} prominent dissociation may be a primary disorder rather than a feature of BPD.
\end{itemize}

\section*{When to Seek Medical Care}
Anyone who recognises persistent instability of mood, identity, or relationships that is causing distress or harming their life should seek a mental health assessment from a GP, psychologist, or psychiatrist. Early engagement with treatment is associated with much better outcomes. A family member or friend who is concerned about a person's self-harm or suicidal behaviour should also encourage and facilitate assessment.

\textbf{Contact your local emergency services for:}
\begin{itemize}
    \item Any thoughts or plans to harm oneself or others
    \item Suicidal behaviour or a suicide attempt
    \item Sudden, severe self-harm needing medical attention
    \item Acute psychotic symptoms, severe confusion, or agitation that puts the person at risk
\end{itemize}
Crisis services such as a local crisis service and the Suicide Call Back Service are available around the clock.

\section*{Diagnosis and Investigations}
\begin{itemize}
    \item \textbf{Clinical interview:} a detailed, non-judgemental assessment of emotions, relationships, self-image, impulsivity, self-harm, and safety, conducted over one or more sessions.
    \item \textbf{Diagnostic criteria:} a formal diagnosis is made when at least five of the specified features are present, are long-standing, and cause significant distress or impairment.
    \item \textbf{Structured interviews and questionnaires:} validated tools such as the Structured Clinical Interview for DSM-5 (SCID) and self-report screens help structure the assessment.
    \item \textbf{Physical health review:} a medical assessment to exclude physical causes of mood and behavioural change, and to document the consequences of self-harm or substance use.
    \item \textbf{Blood tests:} not diagnostic for BPD, but useful to screen for medical conditions such as thyroid disease or anaemia, and for the effects of alcohol or drug use.
    \item \textbf{Assessment of co-morbidity:} screening for depression, anxiety, substance use, and eating disorders, since these are common and affect treatment.
\end{itemize}

\section*{Management}
\begin{itemize}
    \item \textbf{Psychotherapy:} the mainstay of treatment. Dialectical behaviour therapy (DBT) focuses on emotion regulation, distress tolerance, and interpersonal skills; mentalisation-based therapy (MBT) helps people understand their own and others' minds; and cognitive behaviour therapy and schema therapy are also effective.
    \item \textbf{Structured programs:} group-based and day programs deliver skills training and support alongside individual therapy.
    \item \textbf{Crisis planning:} a written plan identifying triggers, early warning signs, coping strategies, and who to contact in a crisis.
    \item \textbf{Self-harm reduction:} teaching safer ways to manage overwhelming emotion and reducing the medical harm of self-injury.
    \item \textbf{Medication:} no drug is specifically licensed for BPD, but medication may be used short-term for co-morbid depression, anxiety, or insomnia, or for severe psychotic-like symptoms.
    \item \textbf{Treatment of co-morbidities:} substance use, eating disorders, and PTSD all need their own treatment, ideally within an integrated plan.
    \item \textbf{Family and carer support:} education and support for families improves outcomes and reduces carer burden.
    \item \textbf{Follow-up:} consistent, long-term engagement is key, since change happens gradually over months to years.
\end{itemize}

\section*{Complications}
\begin{itemize}
    \item Self-harm and suicide; up to about 10\% of people with BPD die by suicide
    \item Impulsive risk-taking with physical injury
    \item Depression, anxiety, and post-traumatic stress disorder
    \item Substance misuse and dependence
    \item Eating disorders and obesity-related illness
    \item Relationship breakdown, isolation, and unemployment
    \item Occupational and financial problems from instability and impulsivity
    \item Repeated hospital admissions and the health consequences of self-harm
\end{itemize}

\section*{Prognosis and Outcomes}
The prognosis of BPD is far better than is commonly believed. Long-term studies show that the majority of people no longer meet diagnostic criteria within ten years, and most experience substantial reductions in self-harm, impulsivity, and hospitalisation. Improvement is strongly associated with receiving structured psychological treatment and with staying engaged over time. Even in the short term, DBT reduces suicidal behaviour and self-harm within months. Some instability of mood and relationships may persist, but it becomes less disabling, and many people achieve stable relationships, work, and a meaningful life. With early help and good support, the outlook is genuinely hopeful.

\section*{Prevention and Self-Care}
\begin{itemize}
    \item Learn and practise emotion-regulation skills, such as naming feelings and using breathing or grounding techniques during distress.
    \item Build a safety and crisis plan with your therapist, and keep it accessible.
    \item Develop a network of trusted supports, and practise asking for help early rather than waiting for a crisis.
    \item Maintain stable routines of sleep, meals, and activity, which reduce emotional volatility.
    \item Avoid alcohol and drugs, which destabilise mood and increase impulsivity.
    \item Reduce exposure to emotionally charged social media and conflict when you are feeling vulnerable.
    \item Attend therapy sessions consistently, even when you feel better.
    \item Address co-occurring conditions, such as anxiety or substance use, since treating them improves the overall picture.
\end{itemize}

\section*{Sources and Further Reading}
The following sources provide reliable, current information on borderline personality disorder and its treatment.
\begin{itemize}
    \item NHS. \textit{Borderline personality disorder.} https://www.nhs.uk/mental-health/conditions/borderline-personality-disorder/
    \item Mayo Clinic. \textit{Borderline personality disorder: symptoms and treatment.} https://www.mayoclinic.org/
    \item Centers for Disease Control and Prevention. \textit{Mental health conditions and resources.} https://www.cdc.gov/mentalhealth/
    \item World Health Organization. \textit{Mental disorders.} https://www.who.int/health-topics/mental-health
    \item NICE. \textit{Borderline personality disorder: recognition and management.} https://www.nice.org.uk/
    \item MedlinePlus. \textit{Borderline personality disorder.} https://medlineplus.gov/
\end{itemize}
